\chapter{Compilar un perfil y leer lo que sale}
\label{ch:compilar}

\section{Los dos subcomandos}

\cod{alexandriac} tiene dos subcomandos. \cod{validate} ejecuta todas las comprobaciones e imprime los
hallazgos. \cod{compile} hace lo mismo y además escribe los artefactos.

\begin{lstlisting}[style=sh]
cd tutorial/manual/examples
../../../bin/alexandriac validate dsl/metering-1.0.0.yaml
\end{lstlisting}
\begin{lstlisting}[style=out]
ok: 1 documents valid
\end{lstlisting}

\begin{lstlisting}[style=sh]
../../../bin/alexandriac compile -o /tmp/metering dsl/metering-1.0.0.yaml
\end{lstlisting}
\begin{lstlisting}[style=out]
/tmp/metering/profiles/odrl/metering/1.0.0/spec.html
/tmp/metering/profiles/odrl/metering/1.0.0/vocab.ttl
/tmp/metering/profiles/odrl/metering/1.0.0/context.jsonld
/tmp/metering/profiles/odrl/metering/1.0.0/prof.ttl
/tmp/metering/profiles/odrl/metering/1.0.0/schema.json
/tmp/metering/profiles/odrl/metering/1.0.0/shapes.ttl
\end{lstlisting}

\noindent\begin{tabularx}{\linewidth}{@{}L{4.6cm}Y@{}}
\toprule
\textbf{Opción} & \textbf{Para qué} \\
\midrule
\texttt{-o <dir>} & Directorio de salida de \cod{compile}. Por defecto, \cod{out}. \\
\texttt{-vendor <dir>} & Copias locales de perfiles externos, para \emph{bindings} vendorizados (capítulo~\ref{ch:spaceprofile}). \\
\texttt{-previous <fichero>} & Fuente de la versión anterior, para comprobar el versionado (capítulo~\ref{ch:evolucion}). Repetible. \\
\cod{-log-level}, \cod{-log-format} & Registro estructurado por la salida de error. \\
\bottomrule
\end{tabularx}

\medskip
El código de salida es \cod{0} si no hay errores, \cod{1} si los hay y \cod{2} si el comando está mal invocado.
Los avisos no cambian el código de salida.

El directorio de un perfil es \texttt{profiles/\allowbreak{}odrl/\allowbreak{}<token>/\allowbreak{}<versión>/\allowbreak{}}, donde el \emph{token} es el último
segmento del \cod{id} del perfil: \cod{metering}.

\section{Los seis artefactos de un perfil}

De 114 líneas de YAML salen unos 37~KB repartidos en seis ficheros. Todos dicen lo mismo en idiomas distintos,
para consumidores distintos.

\noindent\begin{tabularx}{\linewidth}{@{}L{2.9cm}L{2.9cm}Y@{}}
\toprule
\textbf{Fichero} & \textbf{Rol} & \textbf{Para quién} \\
\midrule
\cod{vocab.ttl} & \cod{vocabulary} & Quien razona sobre el vocabulario: ontología OWL, colecciones SKOS y firmas. \\
\cod{context.jsonld} & \cod{mapping} & Quien escribe mensajes: un alias corto por término. \\
\cod{schema.json} & \cod{json-schema} & El gate que valida la forma compacta de una política. \\
\cod{shapes.ttl} & \cod{shacl} & El gate que validaría el grafo expandido (capítulo~\ref{ch:shacl}). \\
\cod{prof.ttl} & \cod{prof} & Catálogos y buscadores: manifiesto W3C PROF del perfil. \\
\cod{spec.html} & \cod{specification} & Personas: una tabla legible de términos. \\
\bottomrule
\end{tabularx}

\subsection{\texttt{vocab.ttl}: la ontología}

Cada término creado se emite en la forma RDF de su tipo. Una acción:

\begin{lstlisting}[style=ttl]
enr:forecastDemand
    a odrl:Action, skos:Concept ;
    odrl:includedIn odrl:use ;
    rdfs:isDefinedBy <https://flexgrid.example/odrl/profile/metering> ;
    rdfs:label "Forecast demand"@en, "Predecir la demanda"@es ;
    skos:definition "To train or run load-forecasting models on the metering data."@en .
\end{lstlisting}

Los términos se agrupan en colecciones SKOS por tipo, lo que permite preguntar «qué operandos izquierdos tiene
este perfil» sin conocer su estructura:

\begin{lstlisting}[style=ttl]
<https://flexgrid.example/odrl/profile/metering#leftOperands>
    a skos:Collection ;
    skos:prefLabel "Constraint Left Operands"@en ;
    skos:member odrl:dateTime, odrl:purpose, enr:minAggregation, enr:processingRegion, enr:retentionDays .
\end{lstlisting}

Y cada firma queda reificada, para que una máquina pueda leer el álgebra de un operando:

\begin{lstlisting}[style=ttl]
<https://flexgrid.example/odrl/profile/metering#sig-minAggregation>
    a sig:Signature ;
    sig:leftOperand enr:minAggregation ;
    sig:operator odrl:gteq, odrl:gt ;
    sig:valueKind sig:Literal ;
    sig:allowReference false ;
    sig:datatype xsd:integer .
\end{lstlisting}

De los términos adoptados solo se dice que son miembros. El perfil no afirma nada sobre IRI que no son suyas.

\subsection{\texttt{context.jsonld}: los alias}

\begin{lstlisting}[style=json]
{"@context":{"@version":1.1,
  "aggregate":"http://www.w3.org/ns/odrl/2/aggregate",
  "anonymize":"http://www.w3.org/ns/odrl/2/anonymize",
  "billing":"enr:billing",
  "dateTime":"http://www.w3.org/ns/odrl/2/dateTime",
  "enr":"https://flexgrid.example/ns/metering#",
  "forecastDemand":"enr:forecastDemand",
  "minAggregation":"enr:minAggregation",
  "purpose":"http://www.w3.org/ns/odrl/2/purpose",
  "regionES":"enr:regionES", … }}
\end{lstlisting}

Hay un alias por término, creado o adoptado. Gracias a él una política escribe \texttt{"action": "forecastDemand"}
en lugar de la IRI completa.

\begin{trampaHabitual}[Dos términos no pueden compartir alias]
Si creas un término \cod{purpose} cuando ya has adoptado \cod{odrl:purpose}, los dos querrían el alias
\cod{purpose} y el contexto no sabría distinguirlos. El compilador lo rechaza:

\begin{lstlisting}[style=out]
error   ALIAS_COLLISION   n2-alias.yaml#1:24 /adopted/left_operands/0 enr:purpose and odrl:purpose share the compact alias "purpose"
\end{lstlisting}
\end{trampaHabitual}

El contexto no redefine \cod{use} ni los doce operadores de ODRL, porque el contexto de DSP ya los define y los
protege.

\subsection{\texttt{schema.json}: el gate que hoy se ejecuta}

Es un JSON Schema Draft 2020-12 \textbf{autocontenido}: no tiene ningún \texttt{\$ref} remoto, así que un gate lo
evalúa sin tocar la red. Valida una política ODRL en la forma compacta que viaja en los mensajes DSP.

Comprueba cuatro cosas. La primera, que la política \textbf{declara el perfil}:

\begin{lstlisting}[style=json]
"profile": {"anyOf": [
  {"enum": ["https://flexgrid.example/odrl/profile/metering",
            "https://flexgrid.example/odrl/profile/metering/1.0.0"]},
  {"type": "array", "contains": {"enum": [ … los mismos dos … ]}}
]}
\end{lstlisting}

La segunda, que las \textbf{acciones} están en la lista cerrada, en cualquiera de sus tres grafías: alias, CURIE
o IRI completa. La tercera, que los \textbf{operandos izquierdos} están en su lista. Y la cuarta, que cada
restricción respeta la \textbf{firma} de su operando, con un condicional por firma:

\begin{lstlisting}[style=json]
{
  "if":   {"required": ["leftOperand"],
           "properties": {"leftOperand": {"enum": ["minAggregation", "enr:minAggregation",
                          "https://flexgrid.example/ns/metering#minAggregation"]}}},
  "then": {"required": ["rightOperand"],
           "not": {"required": ["rightOperandReference"]},
           "properties": {
             "operator": {"enum": ["gteq", "odrl:gteq", "http://www.w3.org/ns/odrl/2/gteq",
                                   "gt",   "odrl:gt",   "http://www.w3.org/ns/odrl/2/gt"]},
             "rightOperand": {"anyOf": [
               {"type": "object", "required": ["@value", "@type"],
                "properties": {"@type": {"enum": ["xsd:integer",
                                         "http://www.w3.org/2001/XMLSchema#integer"]},
                               "@value": {"type": "string", "pattern": "^[+-]?[0-9]+$"}}},
               {"type": "integer"}]}}}
}
\end{lstlisting}

Se lee así: \emph{si} el operando izquierdo es \cod{minAggregation}, \emph{entonces} el operador es \cod{gteq} o
\cod{gt}, el operando derecho es obligatorio y es un entero, y no se admite una referencia. Si el entero va
como valor tipado, su \cod{@value} tiene que tener además la forma léxica de un entero: el patrón rechaza
\texttt{"abc"} aunque diga ser \cod{xsd:integer}.

Fuera de las firmas, toda regla tiene que llevar \cod{action} cuando la lista de acciones es cerrada, y toda
restricción atómica tiene que llevar \cod{leftOperand}, \cod{operator} y un operando derecho o una referencia.

\subsection{\texttt{shapes.ttl}: lo mismo, sobre el grafo}

El mismo perfil en SHACL. La firma anterior es esta \emph{shape}:

\begin{lstlisting}[style=ttl]
gen:Sig_minAggregation a sh:NodeShape ;
  sh:targetSubjectsOf odrl:leftOperand ;
  sh:message "Constraint enr:minAggregation: signature not respected." ;
  sh:or (
    [ sh:not [ sh:path odrl:leftOperand ; sh:hasValue enr:minAggregation ] ]
    [
      sh:property [ sh:path odrl:operator ; sh:minCount 1 ; sh:maxCount 1 ; sh:in ( odrl:gteq odrl:gt ) ] ;
      sh:property [ sh:path odrl:rightOperand ; sh:minCount 1 ; sh:maxCount 1 ;
                    sh:nodeKind sh:Literal ; sh:datatype xsd:integer ] ;
      sh:property [ sh:path odrl:rightOperandReference ; sh:maxCount 0 ]
    ]
  ) .
\end{lstlisting}

Es el mismo condicional, escrito como «o no es este operando, o cumple estas tres propiedades». Junto a las
\emph{shapes} de firma, el fichero lleva \cod{gen:ProfileShape}, \cod{gen:ActionShape},
\cod{gen:LeftOperandShape} y \cod{gen:ConstraintShape}, que repiten las demás comprobaciones del esquema. El
capítulo~\ref{ch:shacl} ejecuta estas \emph{shapes} con un motor de verdad y compara los resultados con los del
esquema.

\subsection{\texttt{prof.ttl} y \texttt{spec.html}}

\cod{prof.ttl} es el manifiesto W3C PROF: dice que el perfil es un perfil de ODRL Core, cuál es su \emph{token}
y qué recurso cumple cada rol. \cod{spec.html} es una página con la tabla de términos, su tipo, su estado y su
definición. No tienen efecto en la validación.

\section{Escribir una política contra el perfil}

Con el perfil compilado, esta es una oferta válida. Es lo que la distribuidora ofrece al agregador:

\begin{lstlisting}[style=json]
{
  "@type": "Offer",
  "@id": "urn:uuid:offer-forecast-01",
  "profile": "https://flexgrid.example/odrl/profile/metering",
  "target": "urn:uuid:dataset-qh-load-madrid-2026",
  "permission": [{
    "action": "forecastDemand",
    "constraint": [
      {"leftOperand": "purpose",          "operator": "eq",      "rightOperand": {"@id": "enr:demandForecast"}},
      {"leftOperand": "minAggregation",   "operator": "gteq",    "rightOperand": 50},
      {"leftOperand": "retentionDays",    "operator": "lteq",    "rightOperand": 90},
      {"leftOperand": "processingRegion", "operator": "isAnyOf",
       "rightOperand": [{"@id": "enr:regionES"}, {"@id": "enr:regionEU"}]},
      {"leftOperand": "dateTime",         "operator": "lt",
       "rightOperand": {"@value": "2027-01-01T00:00:00Z", "@type": "xsd:dateTime"}}
    ]
  }]
}
\end{lstlisting}

En castellano: se permite prever la demanda, solo con esa finalidad, sobre registros que agreguen al menos 50
puntos de suministro, guardándolos como mucho 90 días, tratándolos en España o en la UE, hasta final de año.

Tres detalles de forma que conviene interiorizar:

\begin{itemize}[leftmargin=*, itemsep=2pt]
  \item Un valor \textbf{IRI} se escribe como objeto \texttt{\{"@id": …\}}. Una cadena suelta es un literal, y la
        firma lo rechaza.
  \item Un \textbf{entero} se escribe como número JSON, o como valor tipado con \cod{@value} y \cod{@type}.
  \item Una \textbf{fecha} tiene que ir tipada. Una cadena suelta con una fecha dentro es \cod{xsd:string}.
\end{itemize}

\section{Qué acepta y qué rechaza}
\label{sec:tabla-ofertas}

La tabla resume qué hace \cod{schema.json} con variaciones de esa oferta. Todas las filas están ejecutadas; el
capítulo~\ref{ch:shacl} repite la tabla añadiendo la columna de SHACL.

\begin{longtable}{@{}L{8.6cm}C{1.6cm}L{5.2cm}@{}}
\toprule
\textbf{Variación} & \textbf{Esquema} & \textbf{Por qué} \\
\midrule
\endhead
La oferta tal cual & \pasa & \\
Acción \cod{odrl:sell} & \falla & No está en la lista cerrada. \\
\cod{minAggregation} con operador \cod{lt} & \falla & La firma solo admite \cod{gteq} y \cod{gt}. \\
\texttt{minAggregation gteq "fifty"} & \falla & No es un entero. \\
\texttt{minAggregation gteq 50.5} & \falla & No es un entero. \\
\texttt{purpose eq enr:marketing} & \falla & Valor fuera de la enumeración. \\
\texttt{purpose eq "enr:billing"}, como cadena & \falla & Tiene que ser \texttt{\{"@id": …\}}. \\
\texttt{processingRegion isAnyOf} con dos regiones & \pasa & Operador de conjunto: admite lista. \\
\texttt{retentionDays lteq [30, 90]} & \falla & Sin operador de conjunto no hay listas. \\
Acción \cod{use}, del núcleo de ODRL & \pasa & Un perfil no puede prohibirla. \\
Operando \cod{odrl:spatial} & \falla & No está en la lista cerrada de operandos. \\
Restricción lógica \cod{and} con dos hojas válidas & \pasa & Las restricciones lógicas se validan por dentro. \\
Restricción lógica con una hoja inválida & \falla & \\
\texttt{dateTime lt "2027-01-01T00:00:00Z"}, sin tipar & \falla & Falta \cod{@type}. \\
\cod{duty} con una acción ajena & \falla & Los deberes anidados también se validan. \\
Sin el campo \cod{profile} & \falla & La política tiene que declarar el perfil. \\
\cod{rightOperandReference} en \cod{minAggregation} & \falla & \cod{allow_reference} es \cod{false}. \\
Una regla sin \cod{action} & \falla & La lista de acciones es cerrada. \\
Fecha tipada con valor \texttt{"mañana"} & \falla & No tiene la forma léxica del tipo. \\
\bottomrule
\end{longtable}

\section{Determinismo}

Las mismas fuentes producen los mismos bytes, siempre. Los prefijos se ordenan, los términos se ordenan por IRI
y el JSON se emite en forma canónica.

No es un detalle estético. El hub publica cada fichero por su resumen SHA-256 y firma el manifiesto que los
lista. Si compilar dos veces diera bytes distintos, dos hubs con la misma fuente publicarían paquetes distintos.

\begin{lstlisting}[style=sh]
../../../bin/alexandriac compile -o /tmp/a dsl/metering-1.0.0.yaml > /dev/null
../../../bin/alexandriac compile -o /tmp/b dsl/metering-1.0.0.yaml > /dev/null
diff -r /tmp/a /tmp/b && echo identicos
\end{lstlisting}

\begin{retos}
\begin{enumerate}[leftmargin=*, itemsep=2pt]
  \item Cambia \cod{rules.actions} a \cod{open}, recompila y busca la definición \cod{Action} en
        \cod{schema.json}. ¿Qué ha pasado con la enumeración? ¿Y con \cod{gen:ActionShape} en \cod{shapes.ttl}?
  \item Abre \cod{spec.html} en un navegador. ¿Qué columna distingue un término adoptado de uno creado?
\end{enumerate}
\end{retos}
